El aislamiento entre organizaciones se garantiza mediante un mecanismo
de filtrado automático basado en dos \textit{managers} de Django:
\texttt{TenantManager} (filtrado, asignado por defecto al atributo
\texttt{objects}) y \texttt{UnfilteredManager} (sin filtro, asignado al
atributo \texttt{unfiltered}). Este mecanismo asegura que la
\textit{query} estándar \texttt{MyModel.objects.all()} devuelve
exclusivamente registros de la organización activa, sin que el código
de las vistas o los servicios deba aplicar el filtro manualmente.

El contexto del \dfn{tenant} se almacena en una variable de hilo y
\texttt{TenantManager} lo consulta perezosamente al evaluar cada
\textit{queryset}. El \cref{code:tenant-managers} muestra la
implementación, casi trivial pero estructuralmente decisiva.

\PythonCode[code:tenant-managers]%
           {\texttt{TenantManager} y \texttt{UnfilteredManager}}%
           {Implementación de los \textit{managers} de aislamiento
            multi-tenant.}%
           {managers.py}{20}{53}{1}

El establecimiento y la limpieza del contexto siguen un flujo que abarca
la autenticación, la ejecución de la vista y el final de la respuesta
(\cref{fig:auth-tenant-flow}):

\begin{enumerate}
  \item \textbf{Llegada de la petición.} El servidor Django ejecuta la
  cadena de \textit{middleware}. En este punto la autenticación de
  \acs{drf} aún no se ha invocado y \texttt{request.user} es
  \texttt{AnonymousUser}; el contexto del \dfn{tenant} permanece vacío.

  \item \textbf{Autenticación.} \acs{drf} invoca
  \texttt{JWTAuthentication.authenticate()}, que extrae el \textit{token}
  Bearer, delega su validación al \dfn{plugin} de autenticación activo,
  carga el usuario con \texttt{User.unfiltered.get(...)} (sin filtro,
  dado que aún no hay contexto) y, si el usuario no es
  superadministrador, establece el contexto de la organización mediante
  \texttt{set\_current\_organization\_id}.

  \item \textbf{Ejecución de la vista.} Cualquier \textit{query}
  realizada por la vista o por los servicios subyacentes a través de
  \texttt{MyModel.objects} queda automáticamente acotada a la
  organización del usuario autenticado. Las vistas no necesitan aplicar
  el filtro de manera explícita.

  \item \textbf{Limpieza.} Al finalizar la respuesta,
  \texttt{ClearTenantContextMiddleware} invoca
  \texttt{clear\_current\_organization\_id} en un bloque
  \texttt{finally}, lo que previene la fuga del contexto entre
  peticiones cuando el servidor reutiliza el mismo hilo de su
  \textit{pool}.
\end{enumerate}

\begin{figure}[Flujo de autenticación y contexto multi-tenant]%
              {fig:auth-tenant-flow}%
              {Flujo de autenticación \acs{jwt} y establecimiento del
               contexto del \dfn{tenant}.}
  \image{}{}{plantuml/auth_secuence}
\end{figure}

Este diseño es deliberadamente defensivo: al asignar
\texttt{TenantManager} como \textit{manager} por defecto, todo
\texttt{MyModel.objects.filter(...)} queda filtrado por organización
sin que el desarrollador deba recordar añadir el filtro manualmente. El
olvido de un \texttt{filter(organization=...)} en una \textit{query}
manual constituiría una fuga de datos entre organizaciones; con esta
arquitectura ese olvido es estructuralmente imposible. El
\textit{manager} \texttt{unfiltered} se reserva para situaciones como
las tareas asíncronas que no disponen de contexto de petición y las
\textit{queries} inter-organización, actualmente solo accesibles al
superadministrador.

El compromiso de esta estrategia se reconoció ya en
\cref{SS:SOTA-MULTITENANT}: el aislamiento es lógico (a nivel de
\textit{queryset} filtrado por código) y no físico (segregación de
esquemas o bases de datos por organización). Esta decisión prioriza la
simplicidad operativa y la eficiencia de recursos sobre las garantías
de aislamiento más fuertes, acordes con un primer escenario de
despliegue donde el número de organizaciones inquilinas es reducido.
La arquitectura admite la migración a un modelo de aislamiento más
segregado (esquemas PostgreSQL separados por organización) sin cambios
en el código de negocio, como línea de trabajo futuro.
